\documentclass[11 pt, a4paper]{article}


\NeedsTeXFormat{LaTeX2e}
\usepackage{layout}

% Extensions
\RequirePackage{geometry} 
\geometry{a4paper,tmargin=1.2cm,bmargin=0.8cm,lmargin=1.5cm,rmargin=0.5cm}

% Paramètres du titre
% *******************
\def\classe{$1^{\text{re}}$-Spécialité mathématiques}
\def\titre{Feuille d'exercices : Fonctions trigonométriques}
\def\site{}

\title{\titre}
%\date{~-07-09-12-}
\date{}
\author{\classe}



% Debut du document
% *****************

\begin{document}
\renewcommand{\labelitemi}{$\star$}
\maketitle

% Debut du texte
% **************

\normalsize

\exo[\textbf{Point image d'un nombre réel}]
\begin{minipage}{10cm}
	\begin{enumerate}
		\item Associer à chaque nombre réel son point image sur le cercle trigonométrique ci-contre.
	\begin{center}
		\renewcommand\arraystretch{1.8} 
		\begin{tabular}{|c|*{6}{C{0.7}|}}
			\hline 
			Nombre réel & $\pi$ & $\dfrac{\pi}{2}$& $-\dfrac{\pi}{2}$ & $4\pi$&$\dfrac{3\pi}{2}$& $-3\pi$ \\ 
			\hline 
			Point image & &  &  & && \\ 
			\hline
		\end{tabular}\\
		\medskip
		\begin{tabular}{|c|*{6}{C{0.7}|}}
			\hline 
			Nombre réel & $-\dfrac{5\pi}{2}$ & $\dfrac{\pi}{6}$& $\dfrac{\pi}{4}$ &  $\dfrac{\pi}{3}$& $-\dfrac{2\pi}{3}$& $\dfrac{5\pi}{6}$   \\ 
			\hline 
			Point image & &  &  &&&  \\ 
			\hline
		\end{tabular}
	\end{center}
		\item Donner $5$ nombres réels, $3$ positifs et $2$ négatifs, qui ont:
		\begin{enumerate}
			\item $D$ comme point image:
			\item $H$ comme point image:
			\item $N$ comme point image:
		\end{enumerate} 
	\end{enumerate}
\end{minipage}
\begin{minipage}{8cm}
	\begin{center}
		\includegraphics[width=.9\linewidth]{pointImage.png}
	\end{center}
\end{minipage}

\exo[\textbf{Point image d'un nombre réel}]
\begin{enumerate}
	\item Ecrire le nombre réel $\dfrac{27\pi}{4}$ sous la forme $x+2k\pi$, où $x \in \of{-\pi}{\pi}$ et $k\in \Z$. 
	\item En déduire un nombre réel qui a le même point image que $\dfrac{27\pi}{4}$ sur le cercle trigonométrique.
\end{enumerate}

\exo[\textbf{Point image d'un nombre réel}]
\begin{enumerate}
	\item Ecrire le nombre réel $-\dfrac{101\pi}{3}$ sous la forme $x+2k\pi$, où $x \in \of{-\pi}{\pi}$ et $k\in \Z$.  
	\item En déduire un nombre réel qui a le même point image que $-\dfrac{101\pi}{3}$ sur le cercle trigonométrique.
\end{enumerate}

\exo[\textbf{cosinus et sinus d'un nombre réel}]
On considère le cercle trigonométrique ci-contre.\\
Les points $M$ et $N$ sont les points images respectifs des nombres réel $x_1$ et $x_2$.\\ 
\begin{minipage}{11cm}
	\begin{enumerate}
		\item 
		\begin{enumerate}
			\item Surligner en rouge une ligne qui mesure $x_1$ ul (unité de longueur).
			\item Marquer en rouge un angle qui mesure $x_1$ rad. 
			\item Lire $\cos(x_1)$ et $\sin(x_1)$ :\\ $\cos(x_1) = .........$ ; $\sin(x_1) = .........$
			\item A l'aide de la calculatrice, déterminer une mesure de l'angle $\widehat{IOM}$, en radian (arrondir à 0,1 rad près) puis en degré (arrondir à 0,1° près).\\
			$\widehat{IOM} \approx .......... \textbf{ rad}$ ; $\widehat{IOM} \approx .......... \textbf{°}$
		\end{enumerate}
		\item 
		\begin{enumerate}
			\item Surligner en vert une ligne qui mesure $x_2$ ul (unité de longueur).
			\item Marquer en rouge un angle qui mesure $x_2$ rad. 
			\item Lire $\cos(x_2)$ et $\sin(x_2)$ :\\ $\cos(x_2) \approx .........$ ; $\sin(x_2) = .........$
			\item A l'aide de la calculatrice, déterminer une mesure de l'angle $\widehat{ION}$, en radian (arrondir à 0,1 rad près) puis en degré (arrondir à 0,1° près).\\
			$\widehat{ION} \approx .......... \textbf{ rad}$ ; $\widehat{ION} \approx .......... \textbf{°}$
		\end{enumerate}
	\end{enumerate}
\end{minipage}
\begin{minipage}{9cm}
	\begin{center}
		\includegraphics[width=0.9\linewidth]{cosEtsin.png}
	\end{center}
\end{minipage}

\begin{enumerate}[start=3]
	\item Placer sur le cercle trigonométrique :
	\begin{enumerate}
		\item $R$, le point image du nombre $x_3$ tel que $\cos(x_3)=\dfrac{1}{5}$ et $\sin(x_3) < 0$.
		\item $S$, le point image du nombre $x_4$ tel que $\sin(x_4)=-\dfrac{3}{5}$ et $\cos(x_4) < 0$.
	\end{enumerate}
\end{enumerate}

\exo[\textbf{cosinus et sinus d'un nombre réel}]
On sait d'un réel $x$ que : $x \in \oo{0}{\pi}$ et $\cos x = \dfrac{5}{13}$.\\
Déterminer la valeur exacte de $\sin x$.
		
\exo[\textbf{cosinus et sinus d'un nombre réel}]
On donne : $\cos(\dfrac{\pi}{5}) = \dfrac{1+\sqrt{5}}{4}$.\\
Démontrer que $\sin(\dfrac{\pi}{5}) = \sqrt{\dfrac{5-\sqrt{5}}{8}}$.

\exo[\textbf{$\clubsuit$ Démonstration}]
Dans un repère orthonormé $(O;I,J)$, on note $\mathscr{C}$ le cercle trigonométrique, $M$ le point associé à $\dfrac{\pi}{3}$ sur le cercle $\mathscr{C}$ et $H$ le projeté orthogonal de $M$ sur l'axe des abscisses.
\begin{enumerate}
	\item Faire une figure.
	\item Quelle est la nature du triangle $OMI$ ?
	\item En déduire la longueur $OH$.
	\item Calculer la valeur exacte de la longueur $MH$.
	\item En déduire les valeurs exactes de $\cos\left(\dfrac{\pi}{3}\right)$ et $\sin\left(\dfrac{\pi}{3}\right)$.
\end{enumerate}



\exo[\textbf{$\clubsuit$ Démonstration}]
Dans un repère orthonormé $(O;I,J)$, on note $\mathscr{C}$ le cercle trigonométrique, $M$ le point associé à $\dfrac{\pi}{4}$ sur le cercle $\mathscr{C}$ et $H$ le projeté orthogonal de $M$ sur l'axe des abscisses.\\
En suivante la même démarche que dans l'exercice précédent, trouver les valeurs exactes de $\sin\left(\dfrac{\pi}{4}\right)$ et de $\cos\left(\dfrac{\pi}{4}\right)$.
%\begin{enumerate}
%	\item Faire une figure.
%	\item Quelle est la nature du triangle $OMI$ ?
%	\item En déduire la longueur $OH$.
%	\item Calculer la valeur exacte de la longueur $MH$.
%	\item En déduire la valeur exacte de $\sin\left(\dfrac{\pi}{4}\right)$.
%\end{enumerate}




\exo[\textbf{Valeurs remarquables}]
\begin{enumerate}
\item Tracer le cercle trigonométrique et placer le point $A$ associé au réel $\dfrac{\pi}{6}$.
\item	Placer le point $B$, symétrique de $A$ par rapport à l'axe des abscisses.
Donner le réel associé à ce point dans l'intervalle $]-\pi ; \pi]$.
\item	Placer le point $C$, symétrique de $A$ par rapport à l'axe des ordonnées.
Donner le réel associé à ce point dans l'intervalle $]-\pi ; \pi]$.
\item Placer le point $D$, symétrique de $A$ par rapport à $O$. Donner les réels associés à ce point dans l'intervalle $[0; 2\pi[$ puis dans l'intervalle $]-\pi ; \pi]$.
\end{enumerate}




\colonnes{
\exo[\textbf{Valeurs remarquables}]
\begin{enumerate}
\item  Résoudre sur $]-\pi ; \pi]$ l'équation $\cos(x) =\dfrac{1}{2}$.
\item Résoudre sur $[0; 2\pi[$ l'équation $\cos(x) =\dfrac{\sqrt{2}}{2}$.
\item  Résoudre sur $]-\pi ; \pi]$ l'équation $\sin(x) =-\dfrac{\sqrt{3}}{2}$.
\item Résoudre sur $[0; 2\pi[$ l'équation $\sin(x) =\dfrac{\sqrt{2}}{2}$.
\end{enumerate}
}{
\psset{unit=2.0cm}
\begin{pspicture}(-1.2,-1.2)(1.2,1.2)
\newrgbcolor{bleu}{0.1 0.05 .5}
\newrgbcolor{prune}{.6 0 .48}
\newrgbcolor{rose}{.95 .8 .9}
\def\pshlabel#1{\footnotesize #1}
\def\psvlabel#1{\footnotesize #1}
\psaxes[linewidth=.75pt,labels=none,ticks=none]{->}(0,0)(-1.2,-1.2)(1.2,1.2)
\psaxes[linewidth=1.0pt,linecolor=red]{->}(0,0)(1,1)
\uput[dl](0,0){\footnotesize{O}}\uput[dr](1,0){\footnotesize{\prune $I$}}\uput[ul](0,1){\footnotesize{\prune $J$}}
\pscircle[linewidth=1.25pt, linecolor=bleu,linestyle=solid](0,0){1} 
\psset{linecolor=prune,linewidth=.5pt,linestyle=dashed,labelsep=4pt}
\rput{0}(0,0){\multido{\n=45+90,\i=1+1}{4}{\cnode*(1;\n){2pt}{A\i}}}
\rput{0}(0,0){\multido{\n=30+30,\i=1+1}{12}{\cnode*(1;\n){2pt}{B\i}}}
\ncline{A1}{A2}  \ncline{A2}{A3}  \ncline{A3}{A4}  \ncline{A4}{A1}
\ncline{B1}{B5}  \ncline{B5}{B7}  \ncline{B7}{B11}  \ncline{B11}{B1}
\ncline{B2}{B4}  \ncline{B4}{B8} \ncline{B8}{B10} \ncline{B10}{B2} 
\nput{10}{A1}{\footnotesize{$\dfrac{\pi}{4}$}}\nput{170}{A2}{\footnotesize{$\dfrac{3\pi}{4}$}}\nput{-170}{A3}{\footnotesize{$-\dfrac{3\pi}{4}$}}\nput{-10}{A4}{\footnotesize{$-\dfrac{\pi}{4}$}}
\nput{0}{B1}{\footnotesize{$\dfrac{\pi}{6}$}}\nput{180}{B5}{\footnotesize{$\dfrac{5\pi}{6}$}}\nput{-180}{B7}{\footnotesize{$-\dfrac{5\pi}{6}$}}\nput{0}{B11}{\footnotesize{$-\dfrac{\pi}{6}$}}
\nput{25}{B2}{\footnotesize{$\dfrac{\pi}{3}$}}\nput{145}{B4}{\footnotesize{$\dfrac{2\pi}{3}$}}\nput{-105}{B8}{\footnotesize{$-\dfrac{2\pi}{3}$}}\nput{-105}{B10}{\footnotesize{$-\dfrac{\pi}{3}$}}
\uput[ur](0,1){\footnotesize{$\dfrac{\pi}{2}$}}\uput[ul](-1,0){\footnotesize{$\pi$}}\uput[dl](0,-1){\footnotesize{$-\dfrac{\pi}{2}$}}
\end{pspicture}
}


%\exo[\textbf{Valeurs remarquables}]
%\begin{enumerate}
%\item  Résoudre sur $]-\pi ; \pi]$ l'inéquation $\cos(x) \geqslant -\dfrac{1}{2}$.
%\item  Résoudre sur $]-\pi ; \pi]$ l'inéquation $\sin(x) < 0$.
%\item  En déduire sur $]-\pi ; \pi]$ les solutions du système d'inéquation suivant : 
%$\ds\left\{\begin{array}{l}  
%\cos(x) \geqslant -\dfrac{1}{2} \\ 
%\sin(x) < 0
%\end{array}
% \right.$
%\end{enumerate}

\newpage

\exo[\textbf{Angles associés}]
\begin{minipage}{11.5cm}
	Sur le cercle trigonométrique ci-contre, on a placé le point image $M$ du réel $\dfrac{3\pi}{8}$.\\
	On donne : $\cos \dfrac{3\pi}{8} \approx 0,38$ et $\sin \dfrac{3\pi}{8} \approx 0,92$.
	\begin{enumerate}
		\item Construire le point image du réel $-\dfrac{3\pi}{8}$ puis compléter :\\ $\cos(-\dfrac{3\pi}{8}) \approx ..............$ et $\sin(-\dfrac{3\pi}{8}) \approx ..............$
		\item Construire le point image du réel $\dfrac{3\pi}{8}+\pi$ puis compléter :\\ $\cos(\dfrac{3\pi}{8}+\pi) \approx ..............$ et $\sin(\dfrac{3\pi}{8}+\pi) \approx ..............$
		\item Construire le point image du réel $\pi-\dfrac{3\pi}{8}$ puis compléter :\\ $\cos(\pi-\dfrac{3\pi}{8}) \approx ..............$ et $\sin(\pi-\dfrac{3\pi}{8}) \approx ..............$
		\item Construire le point image du réel $\dfrac{3\pi}{8}+\dfrac{\pi}{2}$ puis compléter :\\ $\cos(\dfrac{3\pi}{8}+\dfrac{\pi}{2}) \approx ..............$ et $\sin(\dfrac{3\pi}{8}+\dfrac{\pi}{2}) \approx ..............$
	\end{enumerate}
\end{minipage}
\begin{minipage}{7cm}
	
	\begin{center}
		\includegraphics[width=0.9\linewidth]{anglesAssociés.png}
	\end{center}
\end{minipage}


\exo[\textbf{Angles associés}]
\begin{enumerate}
	\item Ecrire sous la forme $\dfrac{a\pi}{b}$ où $a \in \Z$ et $b \in \Z, b \neq 0$ :
	\begin{tabenum}
		\item $\pi-\dfrac{2\pi}{5}$ \quad \quad 
		\item $\dfrac{\pi}{2}-\dfrac{2\pi}{5}$ \quad \quad 
		\item $2\pi-\dfrac{2\pi}{5}$ \quad \quad 
		\item $\dfrac{2\pi}{5} + 4\pi$ \quad \quad
		\item $\pi+\dfrac{2\pi}{5}$ \quad \quad  
		\item $\dfrac{\pi}{2}+\dfrac{2\pi}{5}$ \quad \quad
	\end{tabenum}
	\item On donne : $\cos \dfrac{2\pi}{5}=\dfrac{\sqrt{5}-1}{4}$.\\
	A l’aide des égalités précédentes, déterminer, en détaillant, les valeurs de :\\
	\begin{tabenum}
		\item $\cos \dfrac{3\pi}{5}$ \quad \quad 
		\item $\cos \dfrac{22\pi}{5}$ \quad \quad 
		\item $\sin \dfrac{\pi}{10}$ \quad \quad 
		\item $\cos \dfrac{7\pi}{5}$ \quad \quad
		\item $\cos \dfrac{8\pi}{5}$ \quad \quad  
		\item $\sin \dfrac{9\pi}{10}$ \quad \quad
	\end{tabenum}
\end{enumerate}
\begin{small}
\Rpp{
\textbf{Fonction paire, fonction impaire}

On considère une fonction $f$ définie sur l'ensemble $D$.
\begin{itemize}
\item[$\bullet$] On dit que la fonction $f$ est \textbf{paire} si $D$ est centré en $0$ et pour tout $x \in D$, on a $f(-x)=f(x)$.
\item[$\bullet$] On dit que la fonction $f$ est \textbf{impaire} si $D$ est centré en $0$ et pour tout $x \in D$, on a $f(-x)=-f(x)$.
\end{itemize}

%\vspace{-0.6cm}
\begin{center}
Interprétation graphique dans un repère orthogonal
\end{center}


\begin{tabularx}{\linewidth}{>{\centering \arraybackslash}X|>{\centering \arraybackslash}X}
\rule[-2ex]{0pt}{6ex}
\textbf{Fonction paire}

La courbe $\mathscr{C}_f$ est symétrique par rapport à l'axe des ordonnées.
&
\textbf{Fonction impaire}

La courbe $\mathscr{C}_f$ est symétrique par rapport à l'origine du repère.
\\
\newrgbcolor{qqwuqq}{0. 0.39215686274509803 0.}
\newrgbcolor{xdxdff}{0.49019607843137253 0.49019607843137253 1.}
\psset{xunit=0.6cm,yunit=0.4cm,algebraic=true,dimen=middle,dotstyle=o,dotsize=5pt 0,linewidth=1.0pt,arrowsize=2pt 2,arrowinset=0.25}
\begin{pspicture*}(-3.04,-2.06)(3.04,5.96)
%\multips(0,-2)(0,1.0){9}{\psline[linestyle=dashed,linecap=1,dash=1.5pt 1.5pt,linewidth=0.4pt,linecolor=lightgray]{c-c}(-3.04,0)(3.04,0)}
%\multips(-3,0)(0.2,0){31}{\psline[linestyle=dashed,linecap=1,dash=1.5pt 1.5pt,linewidth=0.4pt,linecolor=lightgray]{c-c}(0,-2.06)(0,5.96)}
\psaxes[labelFontSize=\scriptstyle,xAxis=true,yAxis=true,labels=none, Dx=0.2,Dy=1.,ticksize=0pt 0,subticks=2]{->}(0,0)(-3.04,-2.06)(3.04,5.96)[x,2.04] [y,5.8]
\psplot[linewidth=1.pt,linecolor=qqwuqq,plotpoints=200]{-3.040000000000003}{3.040000000000001}{-x^(2.0)+5.0}
\psline[linewidth=1.pt,linestyle=dashed,dash=1pt 1pt](-2.,1.)(2.,1.)
\psline[linewidth=1.pt,linestyle=dashed,dash=1pt 1pt](-2.,0.)(-2.,1.)
\psline[linewidth=1.pt](2.,0.)(2.,1.)
\begin{scriptsize}
\rput[bl](-2.501442622950823,-2.0028774928774937){\qqwuqq{$\mathscr{C}_f$}}
\psdots[dotsize=3pt 0, dotstyle=*,linecolor=xdxdff](-2.,1.)
\rput[bl](-1.2,1.116011396011395){\xdxdff{$f(-x)=f(x)$}}
\psdots[dotsize=3pt 0, dotstyle=*,linecolor=xdxdff](2.,1.)
%\rput[bl](2.0183606557377045,1.116011396011395){\xdxdff{$B$}}
%\rput[bl](0.,0.8189743589743581){$g$}
%\psdots[dotstyle=*,linecolor=xdxdff](-2.,0.)
\rput[bl](-1.9784918032786916,-0.5527065526982){{$-x$}}
%\rput[bl](-1.923672131147544,0.5105128205128197){$h$}
%\psdots[dotstyle=*,linecolor=xdxdff](2.,0.)
\rput[bl](2.0183606557377045,-0.5527065526982){{$x$}}
\rput[bl](-0.3,-0.4){$O$}
%\rput[bl](2.0781639344262293,0.5105128205128197){$i$}
\end{scriptsize}
\end{pspicture*} 
&
\newrgbcolor{wwwwww}{0.4 0.4 0.4}
\newrgbcolor{xdxdff}{0.49019607843137253 0.49019607843137253 1.}
\newrgbcolor{ududff}{0.30196078431372547 0.30196078431372547 1.}
\psset{xunit=0.9cm,yunit=0.7cm,algebraic=true,dimen=middle,dotstyle=o,dotsize=5pt 0,linewidth=1.0pt,arrowsize=2pt 2,arrowinset=0.25}
\begin{pspicture*}(-2.98,-3.)(3.02,3.)
%\multips(0,-3)(0,1.0){7}{\psline[linestyle=dashed,linecap=1,dash=1.5pt 1.5pt,linewidth=0.4pt,linecolor=lightgray]{c-c}(-2.98,0)(3.02,0)}
%\multips(-2,0)(1.0,0){7}{\psline[linestyle=dashed,linecap=1,dash=1.5pt 1.5pt,linewidth=0.4pt,linecolor=lightgray]{c-c}(0,-3.)(0,3.)}
\psaxes[labelFontSize=\scriptstyle,xAxis=true,yAxis=true, labels=none,Dx=1.,Dy=1.,ticksize=0pt 0,subticks=2]{->}(0,0)(-2.98,-3.)(3.02,3.)
\psplot[linewidth=1.pt,linecolor=wwwwww,plotpoints=200]{-2.9800000000000013}{3.020000000000001}{-x^(3.0)+3.0*x}
\psline[linewidth=1.pt](-2.,2.)(2.,-2.)
\psline[linewidth=0.4pt,linestyle=dashed,dash=1pt 1pt](-2.,2.)(-2.,0.)
\psline[linewidth=0.4pt,linestyle=dashed,dash=1pt 1pt](2.,-2.)(0.,-2.)
\psline[linewidth=0.4pt,linestyle=dashed,dash=1pt 1pt](2.,-2.)(2.,0.)
\psline[linewidth=0.4pt,linestyle=dashed,dash=1pt 1pt](-2.,2.)(0.,2.)
\begin{scriptsize}
\rput[bl](2.18,-3.04){\wwwwww{$\mathscr{C}_f$}}
\psdots[dotsize=3pt 0, dotstyle=*,linecolor=xdxdff](-2.,2.)
\rput[bl](-1.92,2.2){\xdxdff{$f(-x)=-f(x)$}}
\psdots[dotsize=3pt 0,dotstyle=*,linecolor=xdxdff](2.,-2.)
%\rput[bl](2.08,-1.8){\xdxdff{$B$}}
%\rput[bl](-0.24,-0.22){$g$}
%\psdots[dotstyle=*,linecolor=xdxdff](-2.,0.)
%\rput[bl](-1.92,0.2){\xdxdff{$C$}}
\rput[bl](-2.32,0.1){$-x$}
%\psdots[dotstyle=*,linecolor=xdxdff](0.,-2.)
\rput[bl](0.08,-2.3){\xdxdff{$f(x)$}}
%\rput[bl](1.,-1.68){$i$}
%\psdots[dotstyle=*,linecolor=ududff](2.,0.)
%\rput[bl](2.08,0.2){\ududff{$E$}}
\rput[bl](2.0,0.1){$x$}
%\rput[bl](-0.2,-0.2){$0$}
\rput[bl](-0.3,-0.3){$O$}
%\psdots[dotstyle=*,linecolor=xdxdff](0.,2.)
%\rput[bl](0.08,2.2){\xdxdff{$F$}}
%\rput[bl](-1.,1.68){$k$}
\end{scriptsize}
\end{pspicture*}
\end{tabularx}
}
\end{small}



%\vspace{-0.9cm}

\newpage
\exo
Soit $f$ la fonction définie sur $\R$ par : $f(x) = \cos(4x)\sin^2(4x)$.

\begin{enumerate}
\item Montrer que $f$ est paire. Interpréter graphiquement.
\item  Montrer que $f$ est $\dfrac{\pi}{2}$-périodique.
\end{enumerate}


\colonnes{
\exo
Soit $g$ la fonction définie sur $\R$ par : \\
\phantom{A} \qquad \qquad $g(x) = \dfrac{1}{4}(3\sin(x) - \sin(3x))$.
\begin{enumerate}
%\item	Tracer la courbe représentative de $g$ sur la calculatrice.
\item	À l'aide de la courbe représentative de $g$ ci-contre, conjecturer la parité et la périodicité de $g$.
\item	Calculer $g(x) + g(-x)$ et en déduire que $g$ est impaire.
\item Montrer que $g$ est périodique et préciser sa période.
\end{enumerate}
}{
\psset{xunit=0.7cm,yunit=0.7cm,algebraic=true,dimen=middle,dotstyle=o,dotsize=5pt 0,linewidth=1.0pt,arrowsize=3pt 2,arrowinset=0.25}
%\psset{unit=0.7cm,arrowsize=7pt}
\begin{pspicture}(-4.0,-1.5)(8.5,1.5)
  \psline{->}(-4.0,0)(8.2,0)
  \psline{->}(0,-1.4)(0,1.4)
  \rput(0.2,-0.3){$O$}
 % plot(\x,{0.25*(3.0*sin(((\x))*180/pi)-sin((3.0*(\x))*180/pi))});
 %  \psplot[plotpoints=1000]{-8}{8}{0.25*(3.0*sin(((\x))*180/pi)-sin((3.0*(\x))*180/pi))}%x 3.14 div 180 mul sin}
%   \draw[line width=2.pt,color=qqwuqq,smooth,samples=100,domain=-4.3:7.3] plot(\x,{0.25*(3.0*sin(((\x))*180/pi)-sin((3.0*(\x))*180/pi))});
  % \psplot[plotpoints=1000]{-4}{8}{x 3.14 div 180 mul sin}
% \psplot[plotpoints=100]{-4}{8}{0.25*(3.0*SIN(x)-SIN(3.0*x))}
 \psplot[linewidth=1.2pt,linecolor=qqwuqq,plotpoints=200]{-4.3}{7.3}{0.25*(3.0*SIN(x)-SIN(3.0*x))}
 \begin{scriptsize}
%  \psline(-6.28,-0.1)(-6.28,0.1)\rput(-6.28,-0.4){$-2\pi$}
 % \psline(-4.59,-0.1)(-4.59,0.1)\rput(-4.59,-0.4){$-\frac{3\pi}{2}$}
  \psline(-3.14,-0.1)(-3.14,0.1)\rput(-3.14,-0.4){$-\pi$}
  \psline(-1.53,-0.1)(-1.53,0.1)\rput(-1.53,-0.4){$-\frac{\pi}{2}$}
  \psline(1.53,-0.1)(1.53,0.1)\rput(1.53,-0.4){$\frac{\pi}{2}$}
  \psline(3.14,-0.1)(3.14,0.1)\rput(3.14,-0.4){$\pi$}
  \psline(4.59,-0.1)(4.59,0.1)\rput(4.59,-0.4){$\frac{3\pi}{2}$}
  \psline(6.28,-0.1)(6.28,0.1)\rput(6.28,-0.4){$2\pi$}
 % \rput(2.8,1.1){$y=\sin x$}
    \psline(-0.1,1.)(0.1,1.0)\rput(-0.28,1.0){\small{$1$}}
    \psline(-0.1,-1.)(0.1,-1.0)\rput(-0.3,-1.0){\small{$-1$}}
    \rput(7.8,-0.3){$x$}
    \end{scriptsize}
\end{pspicture}
}

\medskip

\colonnes{
\exo
On considère la fonction tangente, notée $\tan$, définie par $\tan(x)=\dfrac{\sin(x)}{\cos(x)}$.
\begin{enumerate}
\item	Résoudre sur $\R$ l'équation $\cos(x) = 0$. En déduire le domaine de définition de la fonction tangente.
\item	Montrer que $\tan$ est impaire. Interpréter graphiquement.
\item	Montrer que tan est $\pi$-périodique. Interpréter graphiquement.
\item	\pen Compléter le tableau ci-contre.
\item Tracer la représentation graphique de $\tan$ sur $\left]\dfrac{-\pi}{2};\dfrac{3\pi}{2}\right[$.
\end{enumerate}
}{
\begin{tabular}{|p{1.5cm}|p{1.cm}|p{1.0cm}|p{1.0cm}|p{1.0cm}|}
%\begin{tabular}{|p{1.5cm}|*{4}{p{1cm}|}
\hline 
\phantom{$\ds\dfrac{A}{A}$} $x$ & \centering $0$ & \centering $\dfrac{\pi}{6}$ & \centering $\dfrac{\pi}{4}$ & ~~ $\dfrac{\pi}{3}$ \\ 
\hline 
\phantom{$\ds\dfrac{A}{A}$} $\cos(x)$ &  &  &  &  \\ 
\hline 
\phantom{$\ds\dfrac{A}{A}$} $\sin(x)$  &  &  &  &  \\ 
\hline 
 \phantom{$\ds\dfrac{1}{1}$} $\tan(x)$  &  &  &  &  \\ 
\hline 
\end{tabular} 
}

\colonnes{
\exo
Soit $f$ la fonction définie sur $\R$ par $f(x) = 2x+ \cos(x)$.\\
On donne la courbe représentative de $f$ ci-contre.
\begin{enumerate}
%\item	À l'aide de la courbe représentative de $f$ ci-contre, conjecturer la parité et la périodicité de $g$.
\item	$f$ est-elle paire ? impaire ? Justifier.
\item	Montrer que, pour tout réel $x$, $2x - 1 \leqslant f(x) \leqslant 2x + 1$. Interpréter graphiquement.
%\item	Écrire un programme en langage Python qui permet de calculer et de stocker dans une liste f[k) pour k entier relatif compris entre-2 et 3.
\end{enumerate}
}{
\newrgbcolor{qqwuqq}{0. 0.39215686274509803 0.}
\psset{xunit=0.6cm,yunit=0.6cm,algebraic=true,dimen=middle,dotstyle=o,dotsize=5pt 0,linewidth=1.0pt,arrowsize=3pt 2,arrowinset=0.25}
\begin{pspicture*}(-2.18,-3.32)(2.38,3.22)
\multips(0,-3)(0,1.0){7}{\psline[linestyle=dashed,linecap=1,dash=1.5pt 1.5pt,linewidth=0.4pt,linecolor=gray]{c-c}(-2.18,0)(2.38,0)}
\multips(-2,0)(1.0,0){5}{\psline[linestyle=dashed,linecap=1,dash=1.5pt 1.5pt,linewidth=0.4pt,linecolor=gray]{c-c}(0,-3.32)(0,3.22)}
\psaxes[labelFontSize=\scriptstyle,xAxis=true,yAxis=true,Dx=1.,Dy=1.,ticksize=-2pt 0,subticks=2]{->}(0,0)(-2.18,-3.32)(2.38,3.22)
\psplot[linewidth=1.2pt,linecolor=qqwuqq,plotpoints=200]{-2.18}{2.3800000000000026}{2.0*x+COS(x)}
\begin{scriptsize}
\rput[bl](-2.08,-5.06){\qqwuqq{$f$}}
\end{scriptsize}
\end{pspicture*}
}
















\renewcommand{\fin}{{\scriptsize \site}\hfil {\scriptsize  \titre}\hfil \upshape {\thepage\slash \pageref{LastPage}}}


\end{document}




